\section{Related Works}
\label{sec:related_works}

Facial demographic attribute recognition has transitioned from handcrafted descriptors toward deep representations. Local region-based approaches leverage the geometry of the central facial region to recognize ethnicity and mitigate background interference \cite{ref11}. However, CNN-based approaches still face challenges in maintaining performance under illumination variations and facial characteristic discrepancies across racial groups, including the Other-Race Effect phenomenon. Increasing the depth of convolutional networks can enrich visual feature representations to differentiate diverse populations \cite{ref13}. In addition, parallel feature-sharing mechanisms with visual attention have been developed to capture inter-attribute correlations while mitigating local representation limits. This approach reflects a paradigm shift from feature engineering toward representation learning that is more adaptive to facial diversity \cite{ref16}.

The application of ViT in facial biometrics offers visual modeling capabilities distinct from convolutional approaches. Through the Multi-Head Self-Attention (MHSA) mechanism, transformers can model inter-patch relationships globally, thereby capturing broader spatial dependencies compared to local convolution operations \cite{ref14}. These characteristics demonstrate the potential of transformers in modeling demographic attributes, including gender, on standardized-resolution facial images \cite{ref17}. Beyond standard ViT architectures, hybrid architectures such as MaxViT combine window-based and grid-based attention to construct hierarchical spatial representations \cite{ref18}. Although transformers possess powerful global representation capabilities, their performance remains influenced by the quality, scale, and distributional characteristics of the training data utilized during the learning process.

Performance disparity investigations in computer vision systems have driven evaluation protocols that account for balanced class and demographic subgroup distributions. Testing on balanced datasets provides controlled experimental conditions for comparing performance across subgroups by mitigating the confounding impact of sample size imbalances. However, distributional balance does not inherently eliminate social biases embedded within the data. Latent representation analysis further demonstrates that transformer models can reproduce social biases originating from large-scale pre-training corpora \cite{ref14}. To reduce performance disparities across racial and gender groups, generative augmentation based on StyleGAN2 has been employed to synthesize variations that help mitigate classification errors \cite{ref15}. These findings underscore the critical importance of systematically conducting evaluations across demographic subgroups.

The limitations of single-domain representations in capturing the diversity of facial characteristics have motivated the exploration of multi-task learning and representation fusion from multiple information sources. Multi-task approaches have been utilized to learn correlations between age and gender attributes through jointly processed facial representations \cite{ref12}. Parallel representation-sharing mechanisms have also been developed to combine diverse affective attributes simultaneously, whereas generative approaches can enrich feature variations across demographic groups \cite{ref16}, \cite{ref15}. Furthermore, the integration of ViT-based latent features uniting facial biometric representations and emotional expressions has demonstrated potential for improving intersectional attribute classification \cite{ref19}. Combining biometric representations with age information also exhibits complementary informational contributions, establishing cross-domain fusion as a compelling approach \cite{ref20}.

Once feature representations are obtained, the choice of downstream classifier governs how high-dimensional features are mapped to the decision space. End-to-end approaches commonly employ a SoftMax layer to map feature representations directly to class probabilities \cite{ref11}, \cite{ref18}. Alternatively, Support Vector Machines (SVM) can learn non-linear decision boundaries within the feature space without updating the frozen backbone parameters \cite{ref13}. Tree-based ensemble models such as RF and probabilistic approaches such as GNB can also serve as classifiers with distinct complexity and class-separation characteristics \cite{ref12}. These differences demonstrate that classifier selection substantially impacts the effectiveness of latent representation utilization, particularly when features originate from multiple domains and exhibit high dimensionality.

Based on the reviewed literature, significant opportunities remain to explore facial demographic attribute modeling through the integration of representations across multiple domains. In contrast to prior studies that predominantly rely on single-domain representations or partial domain combinations, this study proposes a tri-domain fusion framework that unifies facial biometric geometric representations, affective expression cues, and biological aging morphological characteristics using pre-trained ViT models with frozen backbones. The latent representations from the three domains are subsequently concatenated and fed into an SVM to learn non-linear decision boundaries. Evaluations are conducted on the balanced DemogPairs dataset to benchmark intersectional race and gender classification while systematically analyzing performance disparities across all evaluated subgroups.
